\begin{abstract}
Recent trends in model merging champion online test-time adaptation (TTA) using unsupervised objectives like entropy minimization to dynamically optimize merging coefficients on unlabeled test streams. In this work, we deconstruct this paradigm and expose a critical methodological flaw: the ``no-data'' strawman, where complex, backpropagation-dependent online adaptation is compared solely against a naive, unoptimized uniform baseline. We show that in real-world deployment, practitioners almost always possess access to a tiny labeled validation set (e.g., 5 to 10 samples per task), which can be leveraged offline to find static, optimal merging coefficients. We propose \textbf{Offline Few-Shot Validation Tuning (OFS-Tune)} and demonstrate that by constraining the merging parameter search space using low-degree polynomials, we can completely reject validation noise and prevent overfitting. Across 30 random seeds under standard visual classification benchmarks, OFS-Tune consistently outperforms online TTA methods while using zero test-time compute. More importantly, we stress-test these paradigms under adversarial stream conditions---such as extreme label shift, temporal task clustering, and ultra-small batch sizes---and show that while online TTA catastrophically collapses due to transductive noise and representation drift, OFS-Tune remains perfectly robust, providing a simple, reliable, and zero-overhead baseline for model merging.
\end{abstract}
